\documentclass{article}
\usepackage{amsmath}
\usepackage{amsfonts}
\usepackage{amssymb}
\usepackage{graphicx}
\usepackage{hyperref}
\usepackage[a4paper, margin=1.0in]{geometry}


\title{Singular Perturbations, Moments and Descriptor Systems: Linear}
\author{Yoonjae Hwang}
\date{30 October 2025}

\begin{document}

\maketitle

\section{Linear Systems - Thinking about explicit solutions and moments}
Consider the following linear time-invariant, SISO system:
\label{sec:linear-systems}
\begin{equation}
\begin{aligned}
\dot{x} &= A x + B u, \\
y & = Cx,
\end{aligned}
\end{equation}

\begin{equation}
\begin{aligned}
\dot{w} &= Sw, \\
u &= Lw,
\end{aligned}
\end{equation}
Given initial conditions \(x(0) = x_0\) and \(w(0) = w_0\), the solution to the system is given by
\begin{equation}
x(t) = e^{At} x_0 + \int_0^t e^{A(t-\tau)} B L e^{S\tau} w_0 d\tau.
\end{equation}
The output of the system is given by
\begin{equation}
y(t) = C e^{At} x_0 + \int_0^t C e^{A(t-\tau)} B L e^{S\tau} w_0 d\tau.
\end{equation}
If $\sigma(A) \cup \sigma(S) = \emptyset$, then there exists a unique solution $\Pi$ to the Sylvester equation
\begin{equation}
A \Pi + B L = \Pi S.
\end{equation}
Using this, we can substitute an expression for \(B L\) into the integral term for $x(t)$:
\begin{equation}
\begin{aligned}
x(t) &= e^{At} x_0 + \int_0^t e^{A(t-\tau)} (\Pi S - A \Pi) e^{S\tau} w_0 d\tau \\
&= e^{At} x_0 + \int_0^t e^{A(t-\tau)} \Pi S e^{S\tau} w_0 d\tau - \int_0^t e^{A(t-\tau)} A \Pi e^{S\tau} w_0 d\tau \\
&= e^{At} x_0 + \left( \Pi e^{St} - e^{At} \Pi \right) w_0 .
\end{aligned}
\end{equation}
Forward Moments are defined to be 
\begin{equation}
M_g = C \Pi
\end{equation}


\section{Problem Setup}
\label{sec:problem-setup}
Consider the following linear singularly perturbed system:
\begin{equation}
\begin{aligned}
\dot{x}_1(t) &= A_{11} x_1(t) + A_{12} x_2(t) + B_1 u(t), \\
\varepsilon \dot{x}_2(t) &= A_{21} x_1(t) + A_{22} x_2(t) + B_2 u(t),
\end{aligned}
\end{equation}
where \(x_1(t) \in \mathbb{R}^{n_1}\), \(x_2(t) \in \mathbb{R}^{n_2}\), \(u(t) \in \mathbb{R}^m\), and \(0 < \epsilon \ll 1\).

You can write this as a descriptor system: 
\begin{equation}
    E_1 \dot{x} = A x + B u
\end{equation}
where
\begin{equation}
    E_1 = \begin{bmatrix}
    I_{n_1} & 0 \\
    0 & \varepsilon I_{n_2}
    \end{bmatrix}, \quad
    A = \begin{bmatrix}
    A_{11} & A_{12} \\
    A_{21} & A_{22}
    \end{bmatrix}, \quad
    B = \begin{bmatrix}
    B_1 \\
    B_2
    \end{bmatrix}, \quad
    x = \begin{bmatrix}
    x_1 \\
    x_2
    \end{bmatrix}.
\end{equation}
The input signal is also a linear singularly perturbed system: 
\begin{equation}
\begin{aligned}
\dot{w_1} = S_{11} w_1 + S_{12}w_2 \\
\varepsilon \dot{w_2} = S_{21} w_1 + S_{22} w_2
\end{aligned}
\end{equation}
where \(w_1(t) \in \mathbb{R}^{r_1}\), \(w_2(t) \in \mathbb{R}^{r_2}\), and \(0 < \epsilon \ll 1\). 
The input to the system is given by \(u(t) = L_1 w_1(t) + L_2 w_2(t)\).

The input signal can also be written as a descriptor system like so:
\begin{equation}
    E_2 \dot{w} = S w
\end{equation}
where
\begin{equation}
    E_2 = \begin{bmatrix}
    I_{r_1} & 0 \\
    0 & \varepsilon I_{r_2}
    \end{bmatrix}, \quad
    S = \begin{bmatrix}
    S_{11} & S_{12} \\
    S_{21} & S_{22}
    \end{bmatrix}, \quad
    w = \begin{bmatrix}
    w_1 \\
    w_2
    \end{bmatrix}.
\end{equation}
The output of the system is given by
\begin{equation}
    y(t) = C_1 x_1(t) + C_2 x_2(t).
\end{equation}

\section{Wong Sequences}
Stuff on Wong Sequences

\section{Wong Sequences: Regular}
We now assume that the pencils $(E_1, A)$ and $(E_2, S)$ are regular. 
Then as shown in the summary by Stephan Trenn, Theorem 2.12: 
\begin{equation}
    P(sE_1-A)Q = s
    \begin{bmatrix}
        I & 0 \\
        0 & N_1
    \end{bmatrix}
    -
    \begin{bmatrix}
        J_1 & 0 \\ 
        0 & I
    \end{bmatrix}.
\end{equation}

\begin{equation}
    W(sE_2-S)T = s
    \begin{bmatrix}
        I & 0 \\
        0 & N_2
    \end{bmatrix}
    -
    \begin{bmatrix}
        J_2 & 0 \\ 
        0 & I
    \end{bmatrix}.
\end{equation}

The following quantities can then be defined for both the system and the signal.
\newline
The index for system is 1, and signal quantities have index 2. 
\newline
The consistency projector is: 
\begin{equation}
    \Xi_1 = Q
    \begin{bmatrix}
        I & 0 \\
        0 & 0
    \end{bmatrix}
    Q^{-1}
\end{equation}

The differential projector:
\begin{equation}
    \Theta_1 = Q
    \begin{bmatrix}
        I & 0 \\
        0 & 0
    \end{bmatrix}
    P
\end{equation}

The consistency projector:
\begin{equation}
    \Phi_1 = Q
    \begin{bmatrix}
        0 & 0 \\
        0 & I
    \end{bmatrix}
    P
\end{equation}

Furthermore: 
\begin{equation}
    \tilde{A} := \Theta_1 A
\end{equation}

\begin{equation}
    \hat{E} := \Phi_1 E_1
\end{equation}

There are some conditions on (P,Q) related to the Wong sequences. 
The explicit solution can then be found as the following: 
Solution for the signal: 
\begin{equation}
    w(t) = e^{\tilde{S} t} \Xi_2 c_2
\end{equation}
The input is then given by: 
\begin{equation}
    f(t) = BLw(t) = BL e^{\tilde{S} t} \Xi_2 c_2
\end{equation}
\begin{equation}
    x(t) = e^{\tilde{A} t} \Xi_1 c + \int_{0}^{t} e^{\tilde{A} (t-\tau)} \Xi_1 BL e^{\tilde{S} \tau} \Xi_2 c_2 d \tau - \sum_{i=0}^{n-1} \hat{E}^i \Phi_1 \left(BL e^{\tilde{S} t} \Xi_2 c_2 \right)^{(i)}
\end{equation}

There are some conditions on $c_1, c_2$ relating the initial conditions. 

We now investigate the integral: 
\begin{equation}
    \int_{0}^{t} e^{\tilde{A} (t-\tau)} \Xi_1 BL e^{\tilde{S} \tau} \Xi_2 c_2 d \tau = e^{\tilde{A}t} \left( \int_{0}^{t} e^{-\tilde{A}\tau} \Xi_1 BL e^{\tilde{S} \tau} d \tau \right) \Xi_2 c_2
\end{equation}

If $\sigma(\tilde{A}) \cup \sigma{\tilde{S}} = \emptyset$:
\begin{equation}
    \tilde{A} \Pi + \Xi_1 BL = \Pi \tilde{S}
\end{equation}

Then the integral evaluates to:
\begin{equation}
    \int_{0}^{t} e^{\tilde{A} (t-\tau)} \Xi_1 BL e^{\tilde{S} \tau} \Xi_2 c_2 d \tau = \left( \Pi e^{\tilde{S} t} - e^{\tilde{A} t} \Pi \right) \Xi_2 c_2
\end{equation}


This is analagous to the result shown in section 1. 


\end{document} 